\begin{textsample}{POS Dim 1 – 2011\_c5 – Score 6.00 – 2011\_c5/enade\_2011\_tec\_red\_comp\_c877d48bcc.txt}  \label{ex:f1_pos_128}
CAPÍTULO 5 DISTRIBUIÇÃO DOS CONCEITOS 5.1 \textbf{PANORAMA} NACIONAL DA DISTRIBUIÇÃO DOS CONCEITOS A Tabela 5.1 apresenta a quantidade e distribuição de cursos da Área de Tecnologia em Redes de Computadores participantes do ENADE/2011, por faixa de conceito e Grande Região.
A diferença entre os cursos \textbf{tabulados} neste capítulo e no capítulo 2 são os cursos sem conceito, em princípio, aqueles sem alunos concluintes que participassem da prova.
Observando -se os dados da Tabela 5.1, nota -se que em todas as regiões, dos 157 cursos participantes, 57 deles ( 36,4 \% ) classificaram -se com conceito 2, o valor modal.
Este foi também o conceito modal nas regiões Sudeste ( 38,6 \%, correspondendo a 29 cursos ) e Centro-Oeste ( 50,0 \%, correspondendo a 12 cursos ).
Em termos nacionais, o conceito 3 ( 35,0 \%, correspondendo a 55 cursos ) foi o segundo mais frequente, seguido dos conceitos 1 e 4 ( 10,8 \%, correspondendo a 17 cursos cada ).
Deste universo de 157 cursos avaliados em Tecnologia em Redes de Computadores, três ( 1,9 \% ) ficaram sem conceito ( SC ).
A região Norte participou com 17 cursos ou 10,8 \% do total nacional.
O conceito modal 3 foi atribuído a sete cursos, o que \textbf{equivale} a 41,1 \% do total regional.
Os demais foram avaliados com conceito 1 ( 17,6 \%, correspondendo a três cursos ), conceito 2 ( 29,4 \%, correspondendo a cinco cursos ) e conceito 5 ( 5,8 \%, correspondendo a um curso ).
Nesta região um curso ( 5,8 \% ) ficou sem conceito.
A região Nordeste participou com 20 cursos ou 12,7 \% do total nacional.
Como já comentado, destes, nove cursos, 45,0 \% em termos regionais, obtiveram conceito 3, o conceito modal para a região.
Os conceitos 2 e 4 foram atribuídos a, respectivamente, seis e 5.2 CONCEITOS POR CATEGORIA ADMINISTRATIVA E POR GRANDE REGIÃO A Tabela 5.2 apresenta a distribuição dos cursos participantes do ENADE/2011 da Área de Tecnologia em Redes de Computadores, por Categoria Administrativa, de acordo com os conceitos por eles alcançados, segundo as Grandes Regiões brasileiras.
Dos 157 cursos participantes 17 ( 10,8 \% ) eram \textbf{ministrados} em instituições públicas e 140 ( 89,2 \% ) em instituições privadas.
Apenas três cursos de IES privadas em todo país não receberam conceito.
De acordo com as informações da Tabela 5.2, em termos nacionais houve uma maior concentração nos conceitos 2 e 3, que somam mais da metade das 157 IES avaliadas.
Dos 17 cursos participantes de IES públicas, o conceito 2 foi o modal, atribuído a oito cursos ( 47,1 \% ).
Entre os demais cursos participantes desta categoria, quatro obtiveram conceito 4 ( 23,5 \% ), três receberam conceito 3 ( 17,6 \% ) e os dois demais receberam conceitos 1 e 5 ( 5,9 \% em cada ).
Na rede privada, o conceito modal foi 3, com 52 cursos ( 37,1 \% ).
Em seguida, o conceito mais frequente foi 2, com 49 cursos ( 35,0 \% ).
O conceito 1 foi atribuído a 16 cursos ( 11,4 \% ), 13 cursos da rede privada receberam o conceito 4 ( 9,3 \% ) e sete cursos alcançaram o conceito 5 ( 5,0 \% ).
Nesta categoria administrativa, três cursos ficaram sem conceito. dois cursos ( 30,0 \% e 10,0 \% ).
Cada um dos conceitos 1 e 5 foram atribuídos a um curso ( 5,0 \% cada um ) da região Nordeste e um curso ficou sem conceito.
Dos 75 cursos participantes da região Sudeste 29 ( 38,6 \% ) obtiveram conceito 2, o conceito modal.
O conceito 3 foi atribuído a 20 cursos ( 26,6 \% ) e o conceito 1 a 11 cursos ( 14,6 \% ).
Dez cursos ( 13,3 \% ) receberam o conceito 4 e cinco outros ( 6,6 \% ) receberam o conceito 5.
Nenhum curso nesta região ficou sem conceito.
A região Sul \textbf{contou} com 21 cursos distribuídos em todas as faixas de conceitos.
A predominância do conceito 3, o valor modal, foi de 42,8 \%, correspondentes a nove participantes na região.
O conceito 2 foi atribuído a cinco cursos ( 23,8 \% ) e o conceito 4, a três cursos ( 14,2 \% ).
O conceito 1 foi recebido por dois cursos ( 9,5 \% ) e o conceito 5 por um curso ( 4,7 \% ).
Um curso ( 4,7 \% ) da região Sul ficou sem conceito.
Dos 24 cursos participantes na região Centro-Oeste, a metade recebeu conceito 2, o conceito modal.
Os demais cursos foram avaliados com conceito 3 ( dez cursos, 41,6 \% ) e conceito 4 ( dois cursos, 8,3 \% ).
Nesta região nenhum curso ficou sem conceito.
Nenhum curso, tampouco, recebeu conceitos 1 ou 5. 5.2 CONCEITOS POR CATEGORIA ADMINISTRATIVA E POR GRANDE REGIÃO A Tabela 5.2 apresenta a distribuição dos cursos participantes do ENADE/2011 da Área de Tecnologia em Redes de Computadores, por Categoria Administrativa, de acordo com os conceitos por eles alcançados, segundo as Grandes Regiões brasileiras.
Dos 157 cursos participantes 17 ( 10,8 \% ) eram \textbf{ministrados} em instituições públicas e 140 ( 89,2 \% ) em instituições privadas.
Apenas três cursos de IES privadas em todo país não receberam conceito.
De acordo com as informações da Tabela 5.2, em termos nacionais houve uma maior concentração nos conceitos 2 e 3, que somam mais da metade das 157 IES avaliadas.
Dos 17 cursos participantes de IES públicas, o conceito 2 foi o modal, atribuído a oito cursos ( 47,1 \% ).
Entre os demais cursos participantes desta categoria, quatro obtiveram conceito 4 ( 23,5 \% ), três receberam conceito 3 ( 17,6 \% ) e os dois demais receberam conceitos 1 e 5 ( 5,9 \% em cada ).
Na rede privada, o conceito modal foi 3, com 52 cursos ( 37,1 \% ).
Em seguida, o conceito mais frequente foi 2, com 49 cursos ( 35,0 \% ).
O conceito 1 foi atribuído a 16 cursos ( 11,4 \% ), 13 cursos da rede privada receberam o conceito 4 ( 9,3 \% ) e sete cursos alcançaram o conceito 5 ( 5,0 \% ).
Nesta categoria administrativa, três cursos ficaram sem conceito.
Na análise por região, observa -se que na região Norte as instituições privadas participaram com 17 cursos e não houve participação de curso vinculado à rede pública.
Desta forma, a distribuição dos conceitos dos cursos de IES Privadas reproduz aquela apresentada para a região como um todo.
Assim, um curso ficou sem conceito ( 5,9 \% ) e os demais receberam conceitos : 1 ( três cursos, 17,6 \% ), 2 ( cinco cursos, 29,4 \% ), 3 ( sete cursos, 41,2 \%, o conceito modal ) e 5 ( um curso, 5,9 \% ). demais receberam conceitos : 1 ( três cursos, 17,6 \% ), 2 ( cinco cursos, 29,4 \% ), 3 ( sete cursos, 41,2 \%, o conceito modal ) e 5 ( um curso, 5,9 \% ).
Na região Nordeste dos 20 cursos participantes dois eram de IES Pública ( 10,0 \% ) e os outros 18 de IES Privada ( 90,0 \% ).
Para as IES Privadas predominaram os cursos com conceito 3, nove cursos, o conceito modal.
Os demais cursos da rede privada receberam os conceitos : 2 ( 33,0 \%, seis curso ), 1 e 4 ( 5,6 \%, correspondendo a um curso cada ).
Um curso desta categoria administrativa ficou sem conceito e nenhum curso recebeu o conceito 5.
Os dois cursos oferecidos por instituições públicas dessa região obtiveram conceitos 4 e 5.
A região Sudeste, como já comentado, foi a que apresentou o maior número de cursos participantes, e a predominância foi também de cursos da rede privada : 66 dos 75 cursos de Tecnologia em Redes de Computadores da região.
Nesta combinação de região e IES Privadas nenhum curso ficou sem conceito e os cursos foram assim classificados : dez cursos com conceito 1, 25 cursos com conceito 2, o conceito modal, 18 com conceito 3, 8 com conceito 4 e 5 com o conceito máximo, 5.
Para as IES \textbf{Públicas} o conceito modal também foi 2, com quatro cursos.
Os demais ficaram assim distribuídos : conceito 1 com um curso e conceitos 3 e 4 com dois cursos cada.
Na região Sul nenhum curso era de IES Pública.
Portanto os 21 cursos da região eram de IES Privadas e a distribuição dos conceitos apresentados na tabela 5.2 reproduz a distribuição dos conceitos da região, apresentados na Tabela 5.1.
Os conceitos foram : nove cursos ( 42,9 \% ) com conceito 3, o modal, cinco cursos ( 23,8 \% ) com conceito 2, três cursos ( 14,3 \% ) no conceito 4, dois cursos ( 9,5 \% ) com conceito 1 por e um com o conceito 5.
Além disso, um dos cursos ficou sem conceito.
Na região Centro-Oeste, 18 dos 24 cursos eram de instituições privadas ( 75,0 \% em termos regionais ) e seis pertencentes a instituições públicas ( 25,0 \% total regional ).
Em instituições privadas, o conceito modal foi 3 obtido pela metade dos cursos desta Categoria Administrativa, ou seja nove cursos.
Dos outros nove cursos um recebeu conceito 4 e oito obtiveram conceito 2.
Dos seis cursos de instituições públicas participantes, quatro foram avaliados com conceito 2 ( 66,7 \% ), o conceito modal.
Um curso recebeu o conceito 3 e outro o conceito 4 ( 16,7 \% cada um ).
Nesta região, nenhum curso de Tecnologia em Redes de Computadores ficou sem conceito ou recebeu o conceito 5. 5.3 CONCEITOS POR ORGANIZAÇÃO ACADÊMICA E POR GRANDE REGIÃO Na Tabela 5.3 encontra -se a distribuição dos conceitos atribuídos aos cursos participantes do ENADE/2011 na Área de Tecnologia em Redes de Computadores, por Organização Acadêmica, segundo as Grandes Regiões brasileiras.
Dos 157 cursos participantes, 42 eram oferecidos em Universidades, 89 em Faculdades e 26 em Centros Universitários.
Esta distribuição corresponde a, respectivamente, 26,8 \%, 56,7 \% e 16,6 \% dos cursos.
Para os 42 cursos das Universidades ( 26,8 \% ) o conceito modal 2 foi obtido por 14 cursos ( 33,3 \% ).
Os demais cursos desta Categoria Administrativa receberam os conceitos : 1 ( 14,3 \%, correspondendo a seis cursos ), 3 ( 26,2 \%, correspondendo a 11 cursos ),, 4 ( 11,9 \%, correspondendo a cinco cursos ) e 5 ( 14,3 \%, correspondendo a seis cursos ).
Não houve curso sem conceito oferecido por Universidade.
Entre os 26 cursos de Centros Universitários o conceito modal também foi 2, com 10 cursos, corresponde a 38,5 \%.
Neste tipo de organização acadêmica um curso ficou sem conceito e nenhum alcançou o conceito 5.
Os demais obtiveram os conceitos : 1(3,8 \%, correspondendo a um curso ), 3 ( 30,8 \%, oito cursos ) e conceito 4(23,1 \%, seis cursos ).
Nas Faculdades, que participaram com o total de 89 cursos ( 56,7 \% ), dois deles ficaram sem conceito.
O conceito 3 foi o modal, com a presença de 36 cursos ( 40,4 \% ).
Os conceitos dos demais cursos neste tipo de Organização Acadêmica foram distribuídos da seguinte forma : 33 receberam o conceito 2 ( 37,1 \% ), dez cursos obtiveram 1(11,2), seis foram avaliados com 4 ( 6,7 \%, cursos ) e dois cursos alcançaram o conceito 5 ( 2,2 \% ).
Considerando -se separadamente as regiões brasileiras, verifica -se que na região Norte apenas dois do total de 17 cursos participantes eram de Universidades.
Um deles recebeu conceito 3 e o outro recebeu o conceito 5.
Os Centros Universitários da região Norte foram representados por quatro cursos ( 23,5 \% ), sendo 3 o conceito modal nesta categoria.
As Faculdades concentraram 11 cursos na região Norte.
O conceito modal foi 2 foi obtido por 5 cursos ( 45,5 \% ) e um curso oferecido em Faculdade na região Norte ficou sem conceito.
Os demais receberam conceito 1 ( três cursos ) e o conceito 3 ( dois cursos ).
Na região Nordeste, as Universidades participaram com quatro cursos na Área de Tecnologia em Redes de Computadores.
Nenhum curso oferecido em Universidades no Nordeste ficou sem conceito.
O conceito modal foi 4, com dois cursos.
Os demais receberam conceito 2 ( um curso ) e conceito 5 ( um curso ).
Os Centros Universitários \textbf{contaram} com dois cursos participantes na região Nordeste, todos eles com conceito 2.
As Faculdades foram representadas por 14 cursos na região Nordeste, para os quais o conceito modal foi 3, com nove cursos.
Dos cursos \textbf{restantes} receberam conceito 2 ( três cursos ) e conceito 1 ( um curso ).
Nesta região houve um curso sem conceito.
Na região Sudeste, as Universidades participaram com 21 dos 75 cursos da região e todos receberam conceito.
Neste grupo seis cursos receberam o conceito modal 1.
Os demais receberam os conceitos 2 ( cinco cursos ), 3 ( quatro cursos ), 4 ( dois cursos ) e 5 ( quatro cursos ).
Os Centros Universitários participaram com 14 cursos na região Sudeste, dos quais sete obtiveram conceito modal, 2, e nenhum ficou sem conceito.
Os demais receberam os conceitos : 4 ( três cursos ), 3 ( três cursos ) e 1 ( um curso ).
As Faculdades foram representadas por 40 cursos na região Sudeste.
Estes se distribuíram nos conceitos : 1 ( quatro cursos ), 2 ( 17 cursos, conceito modal ), 3 ( 13 cursos ), 4 ( cinco cursos ) e o 5 ( um curso ).
Nenhum curso ficou sem conceito.
Dos 21 cursos da região Sul, três eram de Universidades, quatro de Centros Universitários e 14 participantes eram de Faculdades.
Os conceitos dos cursos oferecidos por Faculdades ficaram assim distribuídos : oito receberam o conceito 3, o modal, dois cursos receberam conceito 1, dois outros o conceito 2, um obteve conceito 4 e um alcançou o conceito 5.
Nenhum curso de Faculdade ficou sem conceito.
Dentre os Centros Universitários da região Sul, que tiveram quatro participantes, dois cursos ficaram com o conceito modal, 4.
Um curso recebeu o conceito 2 e outro curso ficou sem conceito.
Já os três cursos oferecidos em Universidades na região obtiveram conceitos 2 ( dois deles ) e 3 ( um curso ).
Na região Centro-Oeste a metade dos 24 cursos participantes eram de Universidades.
O conceito modal foi 2, com seis cursos.
O conceito 3 foi obtidos por cinco cursos e o conceito 4 por apenas um curso.
Nesse tipo de organização, nenhum curso ficou sem conceito.
Os Centros Universitários da região Centro-Oeste \textbf{contaram} com dois cursos, dos quais um recebeu o conceito 3 e o outro o conceito 4.
Os dez cursos participantes oferecidos em Faculdades na região Centro-Oeste se distribuíram em apenas dois conceitos : 2 ( seis cursos, conceito modal ) e 3 ( quatro cursos ).

% matched lemmas: contar, equivaler, ministrar, panorama, pública, restante, tabular
\end{textsample}
